This monograph presents \textsc{mathlib5}, an architecture for verified
symbolic computation organized as a \emph{Verified Symbolic Compute Pipeline}
(VSCP). The central thesis is that confidence in mathematical software is
increased when every transformation between symbolic notation and executable
artifact remains inspectable, provable, and cryptographically attributable.

Traditional software engineering relies on testing, code review, and
operational monitoring to establish confidence. \textsc{mathlib5} explores a
complementary direction: it treats formal verification and provenance as
first-class engineering outputs rather than afterthoughts. The pipeline
separates concerns into ten explicit \emph{safety boundaries}. Instead of a
monolithic compiler, each stage verifies its assumptions before passing
artifacts forward, reducing hidden state and producing a chain of
independently auditable transformations.

The ten layers are: (1) a symbolic \textsc{apl} frontend; (2) a canonical
symbolic intermediate representation; (3) refinement validation with Liquid
Haskell; (4) Lean~4 proof obligations; (5) a verified \textsc{c{--}}
execution layer; (6) logic validation and static analysis; (7) answer-set
programming for policy reasoning; (8) optimization, search, and a P/NP
solver swarm; (9) write-once-read-many (WORM) provenance receipts; and (10)
automated proof-completeness checking.

We describe the mathematical model in which \textsc{apl} acts as a concise
notation for array mathematics, a canonical IR normalizes expressions prior
to proof, refinement types encode structural constraints, and theorem
provers discharge proof obligations before execution artifacts are emitted.
We detail the verified kernel primitives (polynomial arithmetic, matrix
multiplication, symbolic differentiation, algebraic decision procedures,
expression normalization) and the foreign-function bridge that cross-validates
\textsc{c{--}} kernels against Lean~4 by FFI. Every transformation emits a
receipt; hash-linked provenance records form an immutable audit history with
Merkle structures and Ed25519 signatures that detect unauthorized
modification.

The exposition is grounded throughout in the open reference implementation
published at \url{https://github.com/SNAPKITTYWEST/mathlib5}, whose sources
are reproduced in the appendices. We close with an evaluation of
FFI$\leftrightarrow$Lean consistency, benchmark results for closed-form
summation theorems, and a discussion of the constellation of sovereign
repositories in which \textsc{mathlib5} participates.

\paragraph{Keywords.} formal verification; symbolic computation; \textsc{apl};
Lean~4; Liquid Haskell; refinement types; \textsc{c{--}}; WORM ledger;
Ed25519; cryptographic provenance; P/NP swarm; safety boundaries.
